\section{\texorpdfstring{条件期望是对信息 \(\sigma\)-代数的投影}{条件期望是对信息 sigma-代数的投影}}
\label{sec:12-Real-Analysis-Support/05-Measure-Integration-and-Conditional-Expectation/05}

训练一个模型预测送达时长：输入是下单时刻、距离、当时的骑手在线数，损失取均方误差，训到收敛。它收敛到哪个函数？

标准答案是条件期望 \(\E[Y\given X]\)。可翻开定义想验算一下，初等公式写的是
\[
\E[Y\given X=x]=\frac{\E[Y\,\ind_{\set{X=x}}]}{P(X=x)},
\]
而 \(X\) 里有连续的下单时刻，任何一个具体的 \(x\) 都满足 \(P(X=x)=0\)。分母是零，分子也是零。

这个 \(0/0\) 说明，那条初等公式是某个特例的算法，不是条件期望的定义。要问的东西——``掌握了这些信息之后，我对 \(Y\) 的最佳猜测是什么''——在连续情形下依然有意义，需要换一种写法把它接住。

\subsection{不问单点，问可识别的区域}

离散情形能算，靠的不是那个分母，而是``在每个能识别出来的事件上取平均''这件事。事件 \(\set{X=x}\) 之所以能拿来取平均，因为它是你能识别的一块。连续情形下单点识别不出来（概率为零），但正概率的事件仍然识别得出来。

于是把``条件于一个点''改成``条件于一族可识别的事件''。第一节说过，可识别事件的集合就是一个 \(\sigma\)-代数。

设概率空间 \((\Omega,\mathcal F,P)\)，你掌握的信息由子 \(\sigma\)-代数 \(\mathcal G\subseteq\mathcal F\) 表示。若你只看得见 \(X\)，那么 \(\mathcal G=\sigma(X)\)，里面装的正是形如 \(\set{X\in B}\) 的全部事件。

\begin{definition}
设 \(\E[\abs{Y}]<\infty\)，\(\mathcal G\subseteq\mathcal F\) 是子 \(\sigma\)-代数。称随机变量 \(Z\) 为 \(Y\) 关于 \(\mathcal G\) 的\textbf{条件期望}，记 \(Z=\E[Y\given\mathcal G]\)，若 \(Z\) 是 \(\mathcal G\)-可测的，且对每个 \(A\in\mathcal G\) 都有
\[
\int_A Z\,dP=\int_A Y\,dP.
\]
\end{definition}

两条要求各管一头。可测性限制答案能看见什么：\(Z\) 只许用 \(\mathcal G\) 里的信息，不许偷看别处。积分等式要求答案在每一块可识别的区域上都把 \(Y\) 的质量匹配上。只匹配全空间是不够的——常数 \(\E[Y]\) 就能做到，可它一点信息也没用。必须对所有 \(A\in\mathcal G\) 都成立，才算在每一块上都对。

这个写法通篇没有除法。它用积分等式取代了除以概率，而积分对零测集天生免疫：零概率的那些块，任何积分都用不到，定义也就不去规定那里的取值。\(0/0\) 的困难消失了，代价是答案只在几乎处处的意义下唯一——两个都满足定义的 \(Z_1,Z_2\) 必然几乎处处相等，条件期望是一个等价类。

条件概率随之得到定义：\(P(B\given\mathcal G)=\E[\ind_B\given\mathcal G]\)。它是随信息变化的随机变量，不是一个固定的数。

\subsection{退回有限分割，看它给出什么}

新定义得跟旧公式对上。设 \(\mathcal G\) 由有限分割 \(B_1,\ldots,B_m\) 生成，每块概率为正。\(\mathcal G\)-可测意味着在每块上取常数，故 \(Z=\sum_j c_j\ind_{B_j}\)。取 \(A=B_k\) 代入积分等式：
\[
c_k\,P(B_k)=\int_{B_k}Y\,dP=\E[Y\,\ind_{B_k}]
\quad\Longrightarrow\quad
c_k=\frac{\E[Y\,\ind_{B_k}]}{P(B_k)}.
\]
正是分层平均。答案在每个无法进一步细分的块上是常数，常数就是块内均值。若某块概率为零，那里的 \(c_j\) 取什么都行——积分等式不去过问——这与几乎处处唯一的说法一致。

回到第一节那张粗细粒度的图：条件期望做的事，就是在每个信息块上把 \(Y\) 抹平成块内平均。块越细，抹得越少，保留的信息越多。日志加一个字段，就是把某些块劈开，条件期望随之变得更贴近 \(Y\)。

对随机变量记 \(\E[Y\given X]\eqdef\E[Y\given\sigma(X)]\)。答案对 \(\sigma(X)\) 可测，在常见的可测空间里这等价于存在可测函数 \(g\) 使 \(\E[Y\given X]=g(X)\)。写 \(\E[Y\given X=x]=g(x)\) 只是挑了 \(g\) 的一个版本，\(g\) 本身仅在 \(P_X\)-几乎处处确定。开头那个训练任务里的模型逼近的就是这个 \(g\)。

\subsection{存在性由上一节的定理供给}

定义写得出来，不等于满足它的随机变量存在。先设 \(Y\ge 0\)，在 \(\mathcal G\) 上定义
\[
\nu(A)=\int_A Y\,dP,\qquad A\in\mathcal G.
\]
非负性使 \(\nu\) 成为测度，可数可加性从积分的可数可加继承过来。若 \(P(A)=0\)，则 \(\nu(A)=0\)，于是 \(\nu\ll P|_{\mathcal G}\)。Radon--Nikodym 定理给出 \(\mathcal G\)-可测的 \(Z=\dfrac{d\nu}{dP|_{\mathcal G}}\)，它对每个 \(A\in\mathcal G\) 满足 \(\int_A Z\,dP=\nu(A)=\int_A Y\,dP\)——两条要求逐字兑现。一般可积的 \(Y\) 拆成正负两部分分别处理即可。

值得停一秒：条件期望是一个 Radon--Nikodym 导数。它不是先猜条件密度再积分算出来的，是``把 \(Y\) 的质量限制在信息 \(\sigma\)-代数上''这个测度相对于 \(P\) 的导数。它因此不依赖密度是否存在、不依赖离散还是连续、不依赖任何参数化。上一节反复说的``相对于谁''，在这里的答案是：相对于被信息截短之后的那个概率测度。

\subsection{在 \texorpdfstring{\(L^2\)}{L2} 里它就是正交投影}

设 \(Y\in L^2\)。所有平方可积且 \(\mathcal G\)-可测的随机变量构成 \(L^2(\Omega,\mathcal F,P)\) 的一个闭子空间，记作 \(L^2(\mathcal G)\)。定义里的积分等式取 \(A\in\mathcal G\)，等价于对一切有界 \(\mathcal G\)-可测的 \(H\) 成立
\[
\E\bigl[(Y-\E[Y\given\mathcal G])\,H\bigr]=0,
\]
因为指示函数的线性组合能逼近这样的 \(H\)。左边是 \(L^2\) 内积 \(\ip{Y-\E[Y\given\mathcal G]}{H}\)。这句话说的是：残差与子空间里的每一个方向都正交。

\begin{center}
\begin{tikzpicture}[>=stealth,scale=1.0]
\draw[gray] (0,0) -- (7,0) -- (8.6,1.7) -- (1.6,1.7) -- cycle;
\node[gray] at (7.4,0.35) {\footnotesize \(L^2(\mathcal G)\)};
\coordinate (O) at (1.9,0.55);
\coordinate (P) at (5.0,1.0);
\coordinate (Y) at (5.0,3.3);
\fill (O) circle (0.05);
\draw[->,very thick] (O) -- (Y);
\draw[->,thick] (O) -- (P);
\draw[dashed,thick] (P) -- (Y);
\draw (5.0,1.0) -- (5.32,1.0) -- (5.32,1.32) -- (5.0,1.32);
\node at (5.0,3.6) {\footnotesize \(Y\)};
\node at (3.3,0.5) {\footnotesize \(\E[Y\given\mathcal G]\)};
\node at (6.35,2.3) {\footnotesize 残差 \(Y-\E[Y\given\mathcal G]\)};
\node at (3.9,-0.55) {\footnotesize 子空间里装的是所有 \(\mathcal G\)-可测的预测};
\end{tikzpicture}
\end{center}

\emph{条件期望是 \(Y\) 在``可由现有信息表达的预测''这个子空间上的垂足，残差垂直于整个子空间。}

正交性立刻给出勾股分解：对任何 \(\mathcal G\)-可测的 \(H\in L^2\)，
\[
\E\bigl[(Y-H)^2\bigr]
=\E\bigl[(Y-\E[Y\given\mathcal G])^2\bigr]
+\E\bigl[(\E[Y\given\mathcal G]-H)^2\bigr].
\]
第一项与 \(H\) 无关，第二项非负且只在 \(H=\E[Y\given\mathcal G]\) 时为零。所以在全体 \(\mathcal G\)-可测预测中，条件期望是均方误差最小的那个——开头那个训练任务的答案，到这里成了一句几何事实。

这张图与线性代数里的那张是同一张，见 \hyperref[sec:03-Linear-Algebra/03-Orthogonality/02]{``投影与最近点：从一条直线走向一般子空间''}。差别只在子空间的内容：最小二乘的子空间由特征的线性组合张成，条件期望的子空间放宽到全体 \(\mathcal G\)-可测函数。线性回归因此是条件期望在一个更小子空间上的投影，两者的关系在 \hyperref[sec:12-Real-Analysis-Support/07-Lp-Spaces-and-Function-Geometry/04]{``条件期望与线性回归是两种投影''}里展开，而 \(L^2\) 之所以支持这套几何，见 \hyperref[sec:12-Real-Analysis-Support/07-Lp-Spaces-and-Function-Geometry/03]{``\(L^2\) 的几何让投影与正交可用''}。

正交性还是残差诊断的依据。模型若真的收敛到条件期望，残差与任何由输入算出来的量都应当不相关——包括输入的平方、交叉项、分桶指示。检查残差是否与某个特征的函数仍有相关，等于检查投影有没有做完；有残余相关，说明子空间里还有没被用上的方向。这条检查不需要任何分布假设，它就是上面那个内积为零的直接推论。

一句边界：这套几何要求二阶矩存在。条件期望本身只需 \(Y\in L^1\)，在那个更大的空间里没有内积，正交性无从谈起，但``最佳预测''的直觉仍然以别的形式保留——比如在 \(L^1\) 中最小化绝对误差的是条件中位数，投影的说法就换了对象。

\subsection{常用性质都是从两条要求读出来的}

线性、保序、条件 Jensen 不等式都直接从定义得到，验证方式一律是``右边的候选是否 \(\mathcal G\)-可测''加``积分等式是否成立''。三条更值得留意。

已知量可以提出：若 \(H\) 有界且 \(\mathcal G\)-可测，则 \(\E[HY\given\mathcal G]=H\,\E[Y\given\mathcal G]\)。可测性不能省——若 \(H\) 携带 \(\mathcal G\) 看不见的随机性，它在条件下就不是常数，提不出来。

独立信息不改变平均：若 \(Y\) 与 \(\mathcal G\) 独立，则 \(\E[Y\given\mathcal G]=\E[Y]\)。投影的说法是 \(Y\) 与子空间垂直，垂足落在原点。

塔式性质：若 \(\mathcal G\subseteq\mathcal H\)，则
\[
\E\bigl[\E[Y\given\mathcal H]\given\mathcal G\bigr]=\E[Y\given\mathcal G].
\]
先投到大子空间再投到小子空间，等于直接投到小的。嵌套是前提；两个互不包含的 \(\sigma\)-代数，交替条件化一般不能换序，也不会停在任何一个上。

条件 Jensen 取 \(\varphi(x)=x^{2}\) 给出 \(\bigl(\E[Y\given\mathcal G]\bigr)^{2}\le\E[Y^{2}\given\mathcal G]\)，取期望即 \(\Var(\E[Y\given\mathcal G])\le\Var(Y)\)。用部分信息做预测会把波动抹平，预测值比原变量更平缓。这正是勾股分解的另一种读法：总方差被劈成``预测解释掉的''与``残差留下的''两块，与 \hyperref[sec:05-Probability/04-Expectation-and-Moments/04]{``条件期望是利用信息的最佳预测''}里的方差分解是同一件事。

\subsection{一条投影把回归、滤波与鞅串在一起}

给定 filtration \((\mathcal F_t)\) 与可积的 \(Y\)，令 \(M_t=\E[Y\given\mathcal F_t]\)。塔式性质直接给出
\[
\E[M_{t+1}\given\mathcal F_t]=M_t,
\]
即 \((M_t)\) 是鞅。含义是：信息不断增加，预测不断更新，但站在今天回看明天的预测，平均下来仍是今天的值——新信息带来的修正没有系统性的偏向。这类过程的性质见 \hyperref[sec:06-Random-Processes/04-Applications/02]{``Martingale：公平游戏与停时''}。

同一个投影换个 \(\mathcal G\) 就换一个应用。\(\mathcal G\) 由特征生成，投影是回归；\(\mathcal G\) 由到当前时刻为止的观测生成，投影是滤波，Kalman 滤波器算的就是这个垂足在线性子空间里的版本；\(\mathcal G\) 由状态与动作生成，投影出现在 Bellman 方程的 \(\E[V(S_{t+1})\given S_t,A_t]\) 里，Markov 假设断言的正是这个 \(\sigma\)-代数已经够用、更早的历史提供不了额外信息。

也正因为如此，有一件事测度论替不了你做。任何子 \(\sigma\)-代数都能写条件期望，数学上一律合法；但若 \(\mathcal G\) 里没有编码与 \(Y\) 真正相关的信息——日志少打了关键字段，状态表示漏掉了历史变量——算出来的投影仍然是那个子空间里的最佳预测，只是那个子空间离 \(Y\) 很远。选择投影到哪里，是建模判断，不是定理。

最后留一个口子。上面说``信息增加，预测更新''，那么 \(t\to\infty\) 时 \(M_t\) 会不会收敛到 \(Y\)？极限与期望能不能交换，这里又一次成为关键，而控制收敛要求的那张包络并不总有。下一节讨论在没有包络时，什么条件仍能拦住质量。
